\subsection{Pflichtenheft (Auszug)}
\label{app:Pflichtenheft}
In folgendem Auszug aus dem Pflichtenheft wird die geplante Umsetzung der im Lastenheft definierten
Anforderungen beschrieben:

\subsubsection{Umsetzung Funktionale Anforderungen}
\label{sec:UmsetzungFunktionaleAnforderungen}

\subsubsection*{Benutzerverwaltung}
\begin{itemize}
    \item Die Anmeldung erfolgt über die Eingabe einer Ticketnummer und eines Passworts 
    \item Die Benutzeroberfläche ist als Login-Maske mit zwei Eingabefeldern und einem Login-Button umgesetzt
    \item Validierung der Eingaben gegen die Datenbank
    \item Zuordnung einer Benutzerrolle (Besucher, \acs{VIP}, Standpersonal)
    \item Anzeige des aktuellen Guthabens nach erfolgreichem Login
    \item Button zum öffnen eines Dialogs zum Aufladen des Guthabens
    \item \acs{QR}-Code wird pro Ticket lokal aus zugrunde liegenden Daten aus dem Model generiert
\end{itemize}

\subsubsection*{Bestellsystem}
\begin{itemize}
    \item Anzeige aller verfügbaren Verkaufsstände in einer Liste
    \item Auswahl eines Standes über Suchfeld oder Tabellenansicht
    \item Anzeige der Produkte inklusive:
    \begin{itemize}
        \item Name
        \item Wartezeit
        \item Lagerbestand
    \end{itemize}
    \item Produkte können über Buttons dem Warenkorb hinzugefügt werden
    \item Positionen im Warenkorb können über Buttons entfernt werden
    \item Der Bestellung kann über ein Eingabefeld ein Sonderwunsch hinzugefügt werden
    \item Über eine Anzeige werden Gesamtpreis und Gesamtwartezeit angezeigt
    \item \acs{VIP}-Besucher können über einen Schalter die Bestellung als priorisiert markieren
    \item Per Button kann die Bestellung aufgegeben oder abgebrochen werden
    \item Nach Abschluss wird das verfügbare Guthaben validiert und aktualisiert
\end{itemize}

\subsubsection*{Bestellverwaltung}
\begin{itemize}
    \item In der Besucherübersicht werden alle laufenden Bestellungen angezeigt
    \item Folgende Felder werden in der Tabellensicht dargestellt:
    \begin{itemize}
        \item Standnummer
        \item Wartezeit
        \item Status
        \item Bestellnummer
    \end{itemize}
    \item Das Standpersonal kann die, dem Stand zugehörigen, Bestellungen in einer Tabelle innerhalb der Verkäufer-Sicht einsehen
    \item Über die Verkäufer-Sicht kann der Status der Bestellungen (Bestellung aufgegeben, In Bearbeitung, Bereit zur Abholung, Ausgegeben) angepasst werden
\end{itemize}

\subsubsection*{Abholprozess}
\begin{itemize}
    \item Der generierte \acs{QR}-Code wird innerhalb der Besucher-Sicht angezeigt
    \item Verifikation der Bestellung durch Scannen des \acs{QR}-Codes
    \item Validierung der Bestellung im System
    \item Statusänderung auf „Ausgegeben“ durch das Standpersonal
\end{itemize}

\subsubsection*{Bestandsverwaltung}
\begin{itemize}
    \item Die Bestände werden automatisch durch das System verwaltet
    \item Reduzierung des Lagerbestands bei Bestellung
    \item Aktualisierung bei Ausgabe
    \item Kennzeichnung von Produkten mit einer Warnung innerhalb er Verkäufer-Sicht als bei Unterschreitung eines Grenzwerts
\end{itemize}

\subsubsection*{Stornierung und Sonderfälle}
\begin{itemize}
    \item Automatische Stornierung bei nicht ausreichendem Bestand
    \item Bestellung kann durch Besucher storniert werden
    \item Automatische Rückerstattung des Guthabens
    \item Stornierung nicht abgeholter Bestellungen kann durch Standpersonal durchgeführt werden
\end{itemize}

\subsubsection{Umsetzung Nichtfunktionale Anforderungen}
\label{sec:UmsetzungNichtfunktionaleAnforderungen}

\subsubsection*{Benutzbarkeit (Usability)}
\begin{itemize}
    \item Klare Strukturierung in Bereiche (Bestellungen, Warenkorb, Benachrichtigungen)
    \item Konsistente Farbgestaltung
    \item Intuitive Bedienung durch Buttons und Tabellen
\end{itemize}

\subsubsection*{Leistungsanforderungen (Performance)}
\begin{itemize}
    \item Datenbankabfragen werden optimiert ausgeführt
    \item Minimierung von Ladezeiten durch direkte Verarbeitung
    \item Echtzeit-Aktualisierung von Bestellstatus und Beständen
\end{itemize}

\subsubsection*{Zuverlässigkeit}
\begin{itemize}
    \item Korrekte Verarbeitung von Bestellungen ohne Datenverlust
    \item Bestell und Bestandsdaten werden über die Datenbank gespeichert und direkt abgerufen
\end{itemize}

\subsubsection*{Sicherheit}
\begin{itemize}
    \item Authentifizierung über Ticketnummer und Passwort
    \item Zugriffsbeschränkung je nach Benutzerrolle
    \item Validierung aller Eingaben
\end{itemize}

\subsubsection*{Wartbarkeit und Erweiterbarkeit}
\begin{itemize}
    \item Modularer Aufbau (Trennung von GUI, Logik und Datenbank)
    \item Erweiterbarkeit für zusätzliche Funktionen (z. B. weitere Rollen oder Zahlungsarten) kann einfach über die Datenbank durchgeführt werden
\end{itemize}